\chapter{Getting Started}
\label{ch:getting_started}

This chapter describes how to flash the firmware, start the cube for the first time and calibrate
it.

\section{Software requirements}

\begin{itemize}
  \item Arduino IDE~2,
  \item Arduino core for the ESP32~\cite{arduino_esp32}, version 3.x (tested with 3.3.10;
        installed through the Boards Manager),
  \item library MPU6050\_light~\cite{mpu6050light} (tested with 1.2.1; installed through the
        Library Manager).
\end{itemize}
The libraries \cmd{BluetoothSerial}, \cmd{Preferences} and \cmd{Wire} are part of the ESP32 core.

\section{Flashing the firmware}

\begin{enumerate}
  \item Open \cmd{firmware/cubilox\_mk1/cubilox\_mk1.ino} in the Arduino IDE. The folder name must
        match the name of the main file, so do not rename it.
  \item Select the board \enquote{ESP32 Dev Module} and the default partition scheme (Bluetooth
        needs the default scheme; the firmware uses about \SI{85}{\percent} of the program memory).
  \item Connect the ESP32 via USB and upload the sketch.
\end{enumerate}
With the Arduino command line tool, the equivalent command is
\begin{lstlisting}[language=bash, numbers=none]
arduino-cli compile --fqbn esp32:esp32:esp32 firmware/cubilox_mk1
arduino-cli upload  --fqbn esp32:esp32:esp32 -p <port> firmware/cubilox_mk1
\end{lstlisting}

\section{Connecting}

Commands can be sent in two ways:
\begin{itemize}
  \item \textbf{USB:} serial monitor of the Arduino IDE, \SI{115200}{baud}, line ending
        \enquote{newline}. The USB cable disturbs the balance, so this is mainly useful on the
        desk.
  \item \textbf{Bluetooth:} pair the device \cmd{Cubilox\_Mk1} and use a serial terminal app, for
        example \enquote{Serial Bluetooth Terminal} on Android. The ESP32 uses Bluetooth Classic
        (serial port profile), which iPhones do not support.
\end{itemize}

\section{First start}

\begin{enumerate}
  \item Charge the battery fully and connect it.
  \item Lay the cube on any side and do not touch it. After a short beep it calibrates the
        gyroscope (about \SI{3}{s}), then two beeps follow, and a long beep signals that it is ready.
  \item Connect via Bluetooth and send \cmd{status}. On the first start, it reports that the edge
        and the vertex are not calibrated.
  \item Optional: run \cmd{motortest} while holding the cube loosely in your hand. All three
        encoder values should be similar (about 380 on the original cube) and the recommended
        signs should match \cmd{config.h}.
  \item \textbf{Calibrate the edge:} send \cmd{edge}. Within \SI{7}{s}, balance the cube by hand on
        an edge that is parallel to the axis of motor X, as still as possible. A long beep starts
        the measurement; hold the cube still for \SI{3}{s}. Two beeps confirm the calibration.
  \item \textbf{Calibrate the vertex:} send \cmd{vertex} and do the same on the vertex. Find the
        point where the cube balances by itself with almost no force from your fingers.
  \item Place the cube on the edge or vertex and hold it gently near the balance point. When it
        detects the pose, it starts by itself (two short beeps).
\end{enumerate}
The calibration is stored permanently. It only has to be repeated if something changes the
centre of mass, for example if the battery is mounted in a different position. If the vertex
balance is restless or the cube drifts away in one direction, a new \cmd{vertex} calibration
usually helps; several attempts may be needed to find the best balance point.

\section{Beep signals}

\begin{table}[htbp]
  \centering
  \caption{Beep signals of the buzzer.}
  \label{tab:beeps}
  \begin{tabularx}{\textwidth}{lX}
    \toprule
    Signal & Meaning \\
    \midrule
    1 short beep at power-on & IMU found, gyroscope calibration follows \\
    2 beeps, then 1 long beep & gyroscope calibrated, cube ready \\
    2 short beeps & balancing started \\
    1 beep (\SI{200}{ms}) & balancing stopped \\
    1 beep, then ticks every second & calibration: balance the cube now \\
    1 long beep (\SI{150}{ms}) during calibration & measurement starts, hold still \\
    2 beeps after calibration & calibration saved \\
    1 very long beep (\SI{800}{ms}) & calibration rejected \\
    3 short beeps (every \SI{30}{s}) & battery low (below \SI{11.1}{V}) \\
    1 beep every \SI{2}{s} when placed on a pose & battery empty (below \SI{10.5}{V}), no start \\
    1-second beep every \SI{2}{s} after power-on & IMU not found \\
    \bottomrule
  \end{tabularx}
\end{table}

\section{Commands}
\label{sec:commands}

\begin{table}[H]
  \centering
  \caption{Commands (USB or Bluetooth, one command per line).}
  \label{tab:commands}
  \begin{tabularx}{\textwidth}{lX}
    \toprule
    Command & Function \\
    \midrule
    \cmd{help} or \cmd{?} & list of commands \\
    \cmd{status} or \cmd{p} & current gains, calibration status, battery voltage, debug mode \\
    \cmd{edge} & calibrate the edge balance point (\cref{sec:calibration}) \\
    \cmd{vertex} & calibrate the vertex balance point \\
    \cmd{<gain> <value>} & set a gain, e.g.\ \cmd{vk1 95} (see \cref{ch:tuning}) \\
    \cmd{<gain>} & show a gain \\
    \cmd{save} & store all gains in flash \\
    \cmd{reset} & delete the stored gains and use the defaults from \cmd{config.h} \\
    \cmd{debug 0} / \cmd{debug 1} & detailed log output off / on \\
    \cmd{motortest} & short pulse on each motor; shows gyroscope reaction, encoder counts and the
      recommended motor and encoder signs \\
    \cmd{vbatcal <V>} & optional: fine-tune the battery measurement to a multimeter reading, e.g.\
      \cmd{vbatcal 12.6} \\
    \bottomrule
  \end{tabularx}
\end{table}
Commands that write to flash or move the motors (\cmd{edge}, \cmd{vertex}, \cmd{save},
\cmd{reset}, \cmd{motortest}, \cmd{vbatcal}) are only accepted while the cube is not balancing.
Changed gains are lost at the next restart unless they are stored with \cmd{save}.
